\documentclass[12pt,a4paper]{article}
\usepackage[utf8]{inputenc}
\usepackage[vietnamese]{babel}
\usepackage{amsmath}
\usepackage{amsfonts}
\usepackage{amssymb}
\usepackage{graphicx}
\usepackage{hyperref}
\usepackage{geometry}
\usepackage{longtable}
\usepackage{booktabs}
\usepackage{tabularx}
\usepackage{enumitem}

\geometry{
    left=3cm,
    right=2cm,
    top=2.5cm,
    bottom=2.5cm
}

\hypersetup{
    colorlinks=true,
    linkcolor=blue,
    filecolor=magenta,      
    urlcolor=cyan,
    pdftitle={Mô tả bài toán - Ứng dụng Bệnh viện},
}

\begin{document}

\begin{titlepage}
    \centering
    \vspace*{2cm}
    
    {\Large \textbf{TRƯỜNG ĐẠI HỌC...}}\\[0.5cm]
    {\large \textbf{KHOA CÔNG NGHỆ THÔNG TIN}}\\[3cm]
    
    {\Huge \textbf{MÔ TẢ BÀI TOÁN \& YÊU CẦU}}\\[1cm]
    {\Large \textbf{ĐỀ TÀI: CHUYỂN ĐỔI SỐ BỆNH VIỆN - ỨNG DỤNG ĐẶT LỊCH VÀ QUẢN LÝ KHÁM CHỮA BỆNH}}\\[3cm]
    
    \begin{flushleft}
    \large
    \textbf{Giới thiệu sinh viên / Nhóm thực hiện:} \\
    - Nguyễn Văn A - 12345678 \\
    - Trần Thị B - 87654321 \\
    \vspace{1cm}
    \textbf{Giảng viên hướng dẫn:} \\
    - TS. Lê Văn C
    \end{flushleft}
    
    \vfill
    
    {\large \textbf{Năm học 2026 - 2027}}
\end{titlepage}

\tableofcontents
\newpage

\section{Mô tả Bài toán}

\subsection{Bối cảnh}
Hiện nay, việc thăm khám và điều trị tại các bệnh viện và phòng khám đang đối mặt với nhiều thách thức trong khâu quản lý và trải nghiệm người bệnh. Tình trạng quá tải, thời gian chờ đợi kéo dài, thủ tục hành chính phức tạp bằng giấy tờ không chỉ gây mệt mỏi cho bệnh nhân mà còn làm giảm hiệu suất làm việc của đội ngũ y bác sĩ. Hồ sơ bệnh án giấy truyền thống dễ thất lạc, khó tra cứu và không hỗ trợ kết nối thông tin giữa các cơ sở y tế. 

Do đó, việc ứng dụng công nghệ thông tin vào quản trị bệnh viện là nhu cầu cấp thiết nhằm hiện đại hóa quy trình, nâng cao chất lượng dịch vụ y tế và hướng tới mô hình "Bệnh viện thông minh".

\subsection{Mục tiêu đề tài}
Đề tài nhằm xây dựng một hệ thống phần mềm toàn diện phục vụ việc "Chuyển đổi số bệnh viện", với các mục tiêu cụ thể:
\begin{itemize}
    \item \textbf{Đối với bệnh nhân:} Cung cấp ứng dụng di động (Mobile App) giúp chủ động đặt lịch khám, giảm thời gian chờ đợi, dễ dàng theo dõi hồ sơ sức khỏe cá nhân và tương tác với bác sĩ.
    \item \textbf{Đối với y bác sĩ:} Cung cấp nền tảng quản lý (Web Admin) giúp trực quan hóa lịch làm việc, quản lý hồ sơ bệnh án điện tử nhanh chóng, giảm tải công việc hành chính.
    \item \textbf{Đối với ban quản lý:} Nắm bắt tình hình hoạt động của bệnh viện thông qua các báo cáo, thống kê thời gian thực, từ đó đưa ra các quyết định điều hành hiệu quả.
\end{itemize}

\subsection{Phạm vi hệ thống}
Hệ thống bao gồm các thành phần chính:
\begin{enumerate}
    \item \textbf{Mobile App (Ứng dụng di động):} Dành cho đối tượng Bệnh nhân (BN). Hỗ trợ trên các nền tảng iOS và Android.
    \item \textbf{Web Admin (Cổng quản trị web):} Dành cho Bác sĩ, Nhân viên y tế và Ban quản trị (Admin).
    \item \textbf{Backend API \& Database:} Xử lý nghiệp vụ lõi, lưu trữ dữ liệu an toàn và đồng bộ hóa giữa các thiết bị.
\end{enumerate}

\newpage

\section{Yêu cầu Chức năng}

Hệ thống được thiết kế hướng tới 3 nhóm người dùng (Actor) chính: Bệnh nhân, Bác sĩ và Quản trị viên (Admin).

\subsection{Nhóm chức năng dành cho Bệnh nhân (Mobile App)}

\begin{longtable}{|p{1cm}|p{4.5cm}|p{9.5cm}|}
\hline
\textbf{STT} & \textbf{Chức năng} & \textbf{Mô tả} \\
\hline
\endhead

1 & Quản lý Tài khoản & Đăng ký, đăng nhập (Email/Số điện thoại), đặt lại mật khẩu. Cập nhật thông tin hồ sơ cá nhân và ảnh đại diện. \\
\hline
2 & Tìm kiếm Cơ sở y tế/Bác sĩ & Tìm kiếm bác sĩ, chuyên khoa dựa trên nhiều tiêu chí (tên, triệu chứng, đánh giá, vị trí). \\
\hline
3 & Đặt lịch khám & Chọn chuyên khoa, chọn bác sĩ, chọn khung thời gian khám phù hợp và điền trước các triệu chứng nghi ngờ. Hệ thống kiểm tra trùng lặp lịch. \\
\hline
4 & Quản lý lịch khám & Xem danh sách lịch khám sắp tới, lịch sử các lần khám trước. Cho phép hủy lịch (có điều kiện thời gian) hoặc dời lịch khám. \\
\hline
5 & Hồ sơ sức khỏe cá nhân & Lưu trữ và xem lại kết quả chẩn đoán, đơn thuốc, chỉ định cận lâm sàng của các lần khám bệnh. \\
\hline
6 & Hệ thống thông báo & Nhận Push Notification nhắc nhở lịch khám sắp tới, thông báo kết quả xét nghiệm, lịch tái khám. \\
\hline
7 & Đánh giá và Phản hồi & Cho phép bệnh nhân để lại đánh giá, nhận xét (Feedback/Rating) cho bác sĩ và dịch vụ sau khi hoàn tất khám. \\
\hline
8 & Tích hợp Thanh toán (Tùy chọn) & Cổng thanh toán trực tuyến (VNPay/MoMo) để tạm ứng viện phí hoặc thanh toán chi phí khám bệnh. \\
\hline
\end{longtable}

\subsection{Nhóm chức năng dành cho Bác sĩ (Web Admin)}

\begin{longtable}{|p{1cm}|p{4.5cm}|p{9.5cm}|}
\hline
\textbf{STT} & \textbf{Chức năng} & \textbf{Mô tả} \\
\hline
\endhead

1 & Quản lý Lịch làm việc & Đăng ký khung giờ làm việc theo tuần/tháng. Xem trực quan lịch trình khám bệnh hàng ngày. \\
\hline
2 & Xét duyệt Lịch khám & Theo dõi danh sách lịch do bệnh nhân đặt, thực hiện xác nhận (Confirm/Accept) hoặc từ chối (Reject/Cancel) kèm lý do. \\
\hline
3 & Quản lý Bệnh án & Truy cập vào hồ sơ sức khỏe điện tử của bệnh nhân đang khám để xem lịch sử bệnh lý. \\
\hline
4 & Khám bệnh \& Kê đơn & Điền kết quả chẩn đoán, triệu chứng lâm sàng. Tạo đơn thuốc điện tử (chọn thuốc từ danh mục, hướng dẫn sử dụng). \\
\hline
5 & Báo cáo \& Thống kê & Xem thống kê tổng số bệnh nhân đã khám, doanh thu cá nhân (nếu có cơ chế chia sẻ doanh thu). \\
\hline
\end{longtable}

\subsection{Nhóm chức năng dành cho Quản trị viên (Web Admin)}

\begin{longtable}{|p{1cm}|p{4.5cm}|p{9.5cm}|}
\hline
\textbf{STT} & \textbf{Chức năng} & \textbf{Mô tả} \\
\hline
\endhead

1 & Dashboard & Trang tổng quan hiển thị các chỉ số chính (KPIs): Số lượng BN, số lượt khám, doanh thu theo thời gian thực dưới dạng biểu đồ. \\
\hline
2 & Quản lý Người dùng & CRUD (Thêm, Xem, Sửa, Xóa/Khóa) tài khoản Bệnh nhân, Bác sĩ, Nhân viên cơ sở y tế. Phân quyền truy cập. \\
\hline
3 & Quản lý Danh mục & Quản lý các từ điển dữ liệu hệ thống: Danh sách chuyên khoa, Danh mục thuốc, Danh mục dịch vụ khám. \\
\hline
4 & Quản lý Đánh giá & Giám sát các đánh giá của bệnh nhân. Duyệt hoặc ẩn các bình luận không phù hợp. \\
\hline
\end{longtable}

\newpage

\section{Yêu cầu Phi chức năng}

\subsection{Yêu cầu về Hiệu năng (Performance)}
\begin{itemize}
    \item \textbf{Thời gian phản hồi:} Hệ thống phải có khả năng phản hồi phần lớn các giao dịch ($> 90\%$) trong thời gian dưới 2 giây. Các thao tác nặng (xuất báo cáo lớn) không quá 10 giây.
    \item \textbf{Khả năng chịu tải (Concurrency):} Hệ thống backend API cần đảm bảo hoạt động ổn định với ít nhất 500 yêu cầu/giây (RPS) vào các giờ cao điểm đặt lịch khám.
    \item \textbf{Tối ưu hiển thị:} Ứng dụng di động cần đảm bảo độ mượt mà (60 FPS), không có độ trễ lớn khi cuộn danh sách bác sĩ hay tải ảnh hồ sơ.
\end{itemize}

\subsection{Yêu cầu về Bảo mật (Security)}
\begin{itemize}
    \item \textbf{Xác thực \& Phân quyền:} Mọi API kết nối từ client tới server đều phải được xác thực thông qua JWT (JSON Web Token). Phân quyền chặt chẽ (RBAC) để chặn truy cập trái phép vào hồ sơ bệnh án.
    \item \textbf{Bảo mật dữ liệu:} Mật khẩu người dùng phải được băm (hashing) sử dụng thuật toán an toàn (như Bcrypt hoặc Argon2) trước khi lưu vào cơ sở dữ liệu.
    \item \textbf{Tính riêng tư y tế:} Dữ liệu y tế (bệnh án, kết quả khám) luôn phải được coi là dữ liệu nhạy cảm cao, phải tuân thủ chuẩn mức an toàn thông tin (có thể ánh xạ các nguyên tắc HIPAA nếu cần).
\end{itemize}

\subsection{Yêu cầu về Tính khả dụng \& Trải nghiệm người dùng (Usability)}
\begin{itemize}
    \item \textbf{UI/UX:} Giao diện trực quan, dễ sử dụng, phù hợp cả với đối tượng người dùng lớn tuổi. Hỗ trợ font chữ dễ nhìn, độ tương phản tốt.
    \item \textbf{Responsive (Web Admin):} Giao diện dành web cho quản trị và bác sĩ cần tương thích đa thiết bị (PC, tablet) để thuận tiện cho bác sĩ di chuyển trong viện.
    \item \textbf{Ngôn ngữ:} Hỗ trợ đa ngôn ngữ (tiếng Việt mặc định, có thể thêm tiếng Anh là Option).
\end{itemize}

\subsection{Yêu cầu về Khả năng Mở rộng (Scalability \& Maintainability)}
\begin{itemize}
    \item \textbf{Kiến trúc:} Các thành phần Frontend, Backend và Database được thiết kế tách rời hoàn toàn với nhau thông qua API, thuận tiện cho việc nâng cấp.
    \item \textbf{Mã nguồn:} Được tuân thủ theo các Design Pattern, cấu trúc thư mục rõ ràng, có quy tắc conventions thống nhất trong nhóm.
    \item \textbf{Tích hợp:} Sẵn sàng cung cấp các endpoint, webhook để trong tương lai hệ thống có thể kết nối với HIS (Hospital Information System) hoặc các LIS, PACS của bệnh viện khác.
\end{itemize}

\end{document}
